\documentclass[11pt]{article}
\usepackage[margin=0.92in]{geometry}
\usepackage{amsmath,amssymb,amsthm,mathtools}
\usepackage{graphicx}
\usepackage[hidelinks]{hyperref}
\usepackage{xcolor}
\usepackage{microtype}

\newtheorem{theorem}{Theorem}
\newtheorem{proposition}{Proposition}
\newcommand{\R}{\mathbb R}
\newcommand{\E}{\mathbb E}
\newcommand{\Pp}{\mathbb P}
\newcommand{\cE}{\mathcal E}
\newcommand{\cR}{\mathcal R}
\DeclareMathOperator{\rank}{rank}

\title{Exponential semiaxes have an unavoidable square-root gap\\
\large A full counterexample to the question in arXiv:1901.06639}
\author{Autonomous mathematical discovery run \texttt{fa\_banach\_001}\\
\texttt{agent\_lane\_02} \; (GPT5.6)}
\date{11 August 2026}

\begin{document}
\maketitle

\begin{center}
\fcolorbox{black}{gray!8}{\parbox{0.92\linewidth}{
\textbf{Status.} Candidate full counterexample, likely valid, pending human
review.  The source's constant-factor comparison is false for every sequence
of semiaxes comparable to a geometric progression.}}
\end{center}

\section{Source question}

Hinrichs, Krieg, Novak, Prochno, and Ullrich study the ellipsoid
\[
 \cE_\sigma=\left\{x\in\ell_2:
       \sum_{j\ge1}\left(\frac{x_j}{\sigma_j}\right)^2\le1\right\}
\]
and the radius left after $n$ independent standard Gaussian measurements,
\[
 \cR_n(\sigma)=\sup\{\|x\|_2:x\in\cE_\sigma,\ G_nx=0\}.
\]
Optimal information has radius $\sigma_{n+1}$.  On PDF page 8 the source asks
whether
\[
             \E\cR_n(\sigma)\asymp \sigma_{n+1}
\]
also holds when $\sigma_j\asymp a^j$, $0<a<1$; its Corollary 9 gives only
$a^n\lesssim\E\cR_n(\sigma)\lesssim n^2a^n$ \cite{HKNPU}.

\begin{figure}[ht]
  \centering
  \includegraphics[width=\textwidth]{figures/open_problem_crop.png}
  \caption{The complete open-question passage and adjacent Corollary 9 on
  source PDF page 8.}
\end{figure}

\section{Proof intuition}

Use only the first $n+1$ coordinates.  The $n$ Gaussian measurements restricted
to these coordinates have a one-dimensional kernel, and this random null line
is already contained in the full kernel.  A unit direction of the line is
uniform on the sphere, hence has the law $g/\|g\|_2$ for an $(n+1)$-dimensional
standard Gaussian $g$.

For geometric semiaxes, reversing coordinates splits the resulting line
radius into two scales.  Its numerator $\|g\|_2$ is of order $\sqrt n$, whereas
the denominator is a geometrically weighted Gaussian square sum with bounded
mean.  A constant-probability intersection of these two elementary events
therefore forces a $\sqrt n$ loss over the optimal radius.  No analysis of the
rest of the infinite-dimensional kernel is needed.

\section{Counterexample theorem}

\begin{theorem}\label{thm:main}
Let $0<a<1$, and let the nonincreasing positive semiaxes satisfy
\[
              c_0a^j\le \sigma_j\le C_0a^j\qquad(j\ge1)
\]
for constants $c_0,C_0>0$.  Then, for every integer $n\ge22$,
\begin{equation}\label{eq:explicit}
 \E\cR_n(\sigma)
 \ge \frac{c_0\sqrt{1-a^2}}{4\sqrt2}\,
        \sqrt{n+1}\,a^{n+1}.
\end{equation}
Consequently
\[
 \frac{\E\cR_n(\sigma)}{\sigma_{n+1}}
 \ge \frac{c_0\sqrt{1-a^2}}{4\sqrt2\,C_0}\sqrt{n+1}
 \longrightarrow\infty.
\]
In particular, the source's proposed comparison
$\E\cR_n(\sigma)\asymp\sigma_{n+1}$ is false for every exponentially
decaying sequence in the stated class.
\end{theorem}

\begin{proof}
Let $A_n$ be the restriction of $G_n$ to its first $n+1$ columns.  Almost
surely $\rank A_n=n$.  Choose a unit vector $v$ spanning $\ker A_n$ and extend
it by zero after coordinate $n+1$.  Then $G_nv=0$.  The radius of
$\cE_\sigma$ on the line $\R v$ is
\[
 L_n(\sigma)
   :=\left(\sum_{j=1}^{n+1}\frac{v_j^2}{\sigma_j^2}\right)^{-1/2},
 \qquad \cR_n(\sigma)\ge L_n(\sigma).
\]

The law of $A_n$ is invariant under right multiplication by every orthogonal
matrix in dimension $n+1$.  Thus its null line is Haar distributed, and we may
write, in distribution,
\[
                 v=\frac{g}{\|g\|_2},
\]
where $g_1,\dots,g_{n+1}$ are independent $N(0,1)$ variables.  Since
$\sigma_j\ge c_0a^j$,
\begin{align*}
 L_n(\sigma)
 &\ge c_0\frac{\|g\|_2}
          {\bigl(\sum_{j=1}^{n+1}a^{-2j}g_j^2\bigr)^{1/2}} \\
 &=c_0a^{n+1}\frac{\|g\|_2}{Q_n^{1/2}},
 \qquad
 Q_n:=\sum_{r=0}^{n}a^{2r}g_{n+1-r}^2.
\end{align*}
Now
\[
          \E Q_n=\sum_{r=0}^{n}a^{2r}\le\frac1{1-a^2},
\]
so Markov's inequality gives
\begin{equation}\label{eq:qevent}
 \Pp\left\{Q_n\le\frac4{1-a^2}\right\}\ge\frac34.
\end{equation}
Also $\|g\|_2^2$ is chi-square with $n+1$ degrees of freedom.  The elementary
chi-square lower-tail estimate
\[
 \Pp\{\|g\|_2^2<(n+1)/2\}\le e^{-(n+1)/16}
\]
implies, when $n\ge22$,
\begin{equation}\label{eq:normevent}
 \Pp\{\|g\|_2^2\ge(n+1)/2\}\ge\frac34.
\end{equation}
No independence between the events in \eqref{eq:qevent} and
\eqref{eq:normevent} is required: their intersection has probability at least
$1/2$.  On that intersection,
\[
 L_n(\sigma)
 \ge c_0a^{n+1}
      \sqrt{\frac{n+1}{2}}\sqrt{\frac{1-a^2}{4}}
 =c_0a^{n+1}\sqrt{\frac{(n+1)(1-a^2)}8}.
\]
Taking expectations and multiplying by the intersection probability proves
\eqref{eq:explicit}.  Finally,
$\sigma_{n+1}\le C_0a^{n+1}$ gives the displayed diverging ratio.
\end{proof}

\section{Sharper witness asymptotics}

For the exact sequence $\sigma_j=a^j$, the random line itself has a clean
limiting scale.  Let $h_0,h_1,\ldots$ be iid standard Gaussians and set
$Q_\infty=\sum_{r\ge0}a^{2r}h_r^2$.  Reordering coordinates in the formula
above gives a coupling under which
\[
 \frac{L_n}{\sqrt n\,a^{n+1}}
 =\frac{(n^{-1}\sum_{r=0}^{n}h_r^2)^{1/2}}
        {(\sum_{r=0}^{n}a^{2r}h_r^2)^{1/2}}
 \longrightarrow Q_\infty^{-1/2}
 \quad\text{almost surely}.
\]
The limit has a finite, strictly positive mean: indeed
$Q_\infty\ge h_0^2+a^2h_1^2$, whose inverse square root is integrable in two
Gaussian dimensions.  Fatou's lemma and $\cR_n\ge L_n$ therefore yield
\[
 \liminf_{n\to\infty}
 \frac{\E\cR_n(a^j)}{\sqrt n\,a^{n+1}}
 \ge \E Q_\infty^{-1/2}>0.
\]
The explicit elementary estimate in Theorem~\ref{thm:main}, rather than this
limit statement, is the primary counterexample.

\section{Verification, limitations, and novelty}

The proof is finite-dimensional until its final inclusion in $\ker G_n$.
It uses only orthogonal invariance, Markov's inequality, and a standard
chi-square tail bound.  A Monte Carlo script included in the packet confirms
that the witness ratio $L_n/(\sqrt n\,a^{n+1})$ stabilizes numerically; this is
a sanity check, not part of the proof.

The theorem answers the source's yes/no comparison completely and negatively.
It does \emph{not} determine the exact order of $\E\cR_n(\sigma)$ between the
new lower bound $\sqrt n\,a^n$ and the source upper bound $n^2a^n$.

A bounded novelty audit on 11 August 2026 searched the run's registry,
solution, attempt, and proof-gap indexes; the exact arXiv id, title, authors,
and open-problem wording; arXiv queries combining random sections, ellipsoids,
and exponential decay; and all 18 works then indexed by OpenAlex as citing
DOI \texttt{10.1090/tran/8502}.  No later answer or occurrence of this
random-line obstruction was found.  Nearby arXiv:2109.14504 concerns primarily
polynomial semiaxes of generalized $\ell_p$-ellipsoids, and the 2022 survey
arXiv:2209.07266 still records only a polynomial-loss exponential upper bound.
Novelty is plausible but remains subject to specialist review.

\paragraph{Human-review recommendation.}
Verify the inclusion of the finite-coordinate null line in the full kernel,
the uniform-sphere law, and the constants in the Markov/chi-square event
intersection.  If those three points pass, classify this as a full negative
answer.

\begin{thebibliography}{9}
\bibitem{HKNPU}
A. Hinrichs, D. Krieg, E. Novak, J. Prochno, and M. Ullrich,
\emph{Random sections of ellipsoids and the power of random information},
Trans. Amer. Math. Soc. \textbf{374} (2021), no.~12, 8691--8713;
arXiv:1901.06639, doi:10.1090/tran/8502.
\end{thebibliography}

\end{document}
